\section{SMORES-EP Case Study: Model, Contact, and Control}
\label{sec:smores}

\subsection{Platform basis}

\smoresep{} was selected because it exercises nearly every architectural concern in a
compact platform: reusable articulated modules, four connector faces, differential
drive, active reconfiguration, discrete mating orientations, and published connector
and topology research. The official platform description identifies four active DOFs
and rubber-tired left/right wheels \citep{modlabSmoresEp}. The EP-Face is genderless and
switchable, with published holding-force and capture characterization
\citep{tosun2016epface}. Self-assembly work divides docking into navigation, pose
adjustment, and final approach \citep{liu2020parallel,liu2023smoresep}.

The committed pack is a research artifact, not a complete hardware specification. This
section therefore labels parameter provenance. Superscripts mean published
(\sourcefact{}), geometry-derived (\cadfact{}), provisional simulation choice
(\simchoice{}), or measured/tested simulation outcome (\measuredfact{}).

\subsection{Pack contents}

The \pathcode{examples/robot_packs/smores_ep} pack contains:

\begin{itemize}
  \item one module type and one genderless self-compatible connector type;
  \item four connector instances named \code{bottom}, \code{pan}, \code{left}, and
        \code{right};
  \item dock and undock capability descriptions;
  \item a default runtime topology view recipe;
  \item a Fusion-derived URDF and a same-ID physics-specific MJCF asset;
  \item five detailed STL visual meshes and simplified collision/contact assets; and
  \item a URDF backend mapping.
\end{itemize}

The module mass is \SI{0.4386368}{kg}\cadfact{}. LEFT and RIGHT wheel joints and PAN are
continuous; TILT is bounded to $\pm\pi/2$. The wheel speed cap is
$\pi/2\,\mathrm{rad/s}$ (\SI{90}{\degree\per\second})\sourcefact{}. The wheel effort limit
\SI{0.04}{N.m}\simchoice{} and PAN/TILT limits \SI{0.10}{N.m}\simchoice{} are
simulation bootstrap values, not identified motor constants.

\subsection{Connector frames}

\begin{table}[htbp]
\centering
\caption{Committed \smoresep{} connector instances. Positions are in parent-link
coordinates.}
\label{tab:smores-connectors}
\begin{tabularx}{\textwidth}{L{0.12\textwidth}L{0.19\textwidth}L{0.28\textwidth}Y}
\toprule
ID & Parent link & Position (m) and RPY (rad) & Docking axis \\
\midrule
\code{bottom} & \code{base\_link} & $(-0.0107416,0,0)$; $(0,0,\pi)$ & $(-1,0,0)$ \\
\code{pan} & \code{front\_wheel\_1} & $(0.005762,0,0)$; $(0,0,0)$ & $(1,0,0)$ \\
\code{left} & \code{left\_wheel\_1} & $(0,0.04405,0)$; $(0,0,\pi/2)$ & $(0,1,0)$ \\
\code{right} & \code{right\_wheel\_1} & $(0,-0.04405,0)$; $(0,0,-\pi/2)$ & $(0,-1,0)$ \\
\bottomrule
\end{tabularx}
\end{table}

The corrected BOTTOM connector is on the rigid rear $-X$ face of the base, not the
underside. Its plane is derived from the dominant rear surface in the Fusion mesh. This
correction is essential for TOP-to-BOTTOM snake connections: an underside frame turns a
nose-to-tail chain into an unintended vertical or kinked structure.

\subsection{EP-face semantic model}

The connector type is genderless and self-compatible. Allowed roll states are
$\{0,\pi/2,\pi,3\pi/2\}$\sourcefact{}. The current capture proxy is cylindrical with
\SI{6}{mm} position tolerance\simchoice{}, \SI{25}{\degree} orientation tolerance, and
\SI{0.05}{m/s} maximum relative speed. The 25-degree value is motivated by published
residual attraction under angular misalignment \citep{tosun2016epface}, but it should
not be read as a calibrated complete capture model.

Connection intent is fixed, undocking is supported, auto-latch is disabled, measured
alignment is selected, and redock cooldown is \SI{0.1}{s}. Listed normal, shear, and
bending loads are explicitly unverified provisional metadata. The current generic
acceptance model uses one radial position bound; it does not reproduce the literature's
separate normal and lateral capture envelopes.

\subsection{Visual and contact models}

Detailed STLs carry visible shape but no textures or material graph. The URDF assigns a
uniform silver material. In MJCF, visual meshes have zero density and do not carry
contact. Contact uses simpler primitives:

\begin{itemize}
  \item two tire cylinders with radius \SI{0.040}{m}\cadfact{} and half-width
        \SI{0.01045}{m}\cadfact{};
  \item axle track \SI{0.0672}{m}\cadfact{};
  \item a small rear support skid\simchoice{} to make the two-wheel chassis statically
        supportable; and
  \item connector-face collision proxies in a separate category, so faces can meet
        without dragging against the ground.
\end{itemize}

Tire/ground contact uses \code{condim=4} and friction tuple
$(0.05,1.0,0.01,10^{-4},10^{-4})$\simchoice{}. The low/high tangential values establish
anisotropy; the backend rotates the tangent frame with the wheel axle. Skid/ground
friction is $(0.08,0.08,0.005,10^{-4},10^{-4})$\simchoice{}. Wheel damping and armature
are \num{0.001} and \num{0.0005}; PAN/TILT damping and armature are \num{0.01} and
\num{0.0002}\simchoice{}. These quantities were tuned for stable first-pass dynamics and
must eventually be identified from hardware.

\subsection{Differential-drive kinematics}

For wheel radius $r$, track $b$, and left/right angular speeds
$\omega_L,\omega_R$, the ideal planar body twist is

\begin{align}
v &= \frac{r}{2}(\omega_R+\omega_L), \\
\dot\psi &= \frac{r}{b}(\omega_R-\omega_L).
\label{eq:diff-drive}
\end{align}

The inverse targets are

\begin{align}
\omega_R^\star &= \frac{v+\frac{b}{2}\dot\psi}{r}, &
\omega_L^\star &= \frac{v-\frac{b}{2}\dot\psi}{r}.
\label{eq:wheel-targets}
\end{align}

Each wheel uses bounded proportional velocity-to-effort feedback,

\begin{equation}
\tau_i = \operatorname{sat}_{\tau_{\max}}
\left[k_v(\omega_i^\star-\omega_i)\right],
\label{eq:wheel-effort}
\end{equation}

with $k_v=0.1\,\mathrm{N,m/(rad/s)}$\simchoice{} and
$\tau_{\max}=0.04\,\mathrm{N,m}$\simchoice{}. PAN and TILT use bounded PD hold,

\begin{equation}
\tau_h = \operatorname{sat}_{\tau_{h,\max}}
\left[k_p(q^\star-q)-k_d\dot q\right],
\label{eq:hold-effort}
\end{equation}

with $k_p=1.0$, $k_d=0.02$, and \SI{0.1}{N.m} cap\simchoice{}. Current feedback runs at
the solver period for stability. A calibrated actuator model should permit a more
hardware-faithful outer-loop frequency.

\subsection{Rigid-assembly wheel allocation}

The final two Driver-to-Snake actions move a connected three-module component. A
desired planar twist is evaluated at every module center. For module $i$ at position
$\mathbf p_i$, reference point $\mathbf p_0$, reference linear velocity $\mathbf v_0$,
and desired yaw rate $\omega$, the rigid velocity target is

\begin{equation}
\mathbf v_i^\star = \mathbf v_0 +
\omega\,\hat{\mathbf z}\times(\mathbf p_i-\mathbf p_0).
\label{eq:rigid-velocity}
\end{equation}

The allocator projects $\mathbf v_i^\star$ onto each module's planar forward axis and
uses \cref{eq:wheel-targets}. The orthogonal projection is a nonholonomic residual. If
the largest requested residual exceeds \SI{0.08}{m/s}\simchoice{}, the whole twist is
scaled. Wheel targets are then scaled again if any exceeds $\pi/2\,\mathrm{rad/s}$.
Both requested and applied residual maxima are retained for diagnosis.

This is a pragmatic controller for a rigid assembly whose module wheel axes are not
all located at the reference center. It is not an optimal traction allocator, and the
reported residual is a command-space diagnostic, not measured tire slip.

\subsection{Two-module physical scenario}

The first physical scenario uses a stationary module's rear BOTTOM connector and a
moving module's PAN/TOP connector. Creation performs all capability, joint-feedback,
and effort-command preflight before staging the pair once. It then:

\begin{enumerate}
  \item settles both free modules for \SI{0.5}{s};
  \item brakes the target and drives the mover forward through wheel effort;
  \item steers from measured connector geometry and continues until a strict
        \SI{1}{mm} near-contact gate is met;
  \item requests the ordinary measured-frame fixed connection;
  \item holds for \SI{1}{s};
  \item waits \SI{80}{ms}\sourcefact{} as a modeled release delay;
  \item removes the weld; and
  \item reverses \SI{0.04}{m} through wheel effort.
\end{enumerate}

No root pose, root twist, or external wrench is written after initial staging. The
approach and separation therefore depend on gravity, contact, inertia, joint motion,
and effort limits. Magnetic attraction before capture is absent; the accepted latch is
an ideal fixed weld.

\subsection{Driver-to-Snake plan}

The seven-module plan follows the connector substitutions reported by
\citet{liu2019distributed}. The initial Driver tree and replacements are listed in
\cref{tab:snake-actions}. The final physical module order is
$1-3-2-4-5-6-7$.

\begin{table}[htbp]
\centering
\caption{Driver-to-Snake connector replacements. TOP maps to \code{pan}; BOTTOM maps to
\code{bottom}.}
\label{tab:snake-actions}
\begin{tabularx}{\textwidth}{L{0.08\textwidth}YY}
\toprule
Action & Release & Replacement \\
\midrule
1 & \code{module\_1/bottom}--\code{module\_2/pan} &
    \code{module\_1/pan}--\code{module\_3/bottom} \\
2 & \code{module\_7/pan}--\code{module\_5/bottom} &
    \code{module\_7/bottom}--\code{module\_6/pan} \\
3 & \code{module\_2/right}--\code{module\_4/left} &
    \code{module\_2/pan}--\code{module\_4/bottom} \\
4 & \code{module\_5/left}--\code{module\_4/right} &
    \code{module\_5/bottom}--\code{module\_4/pan} \\
\bottomrule
\end{tabularx}
\end{table}

Actions 1 and 2 move single leaves. After they dock, action 3 moves the connected
$\{1,3,2\}$ component, and action 4 moves $\{5,6,7\}$. The plan is separated from both
the kinematic and physical scenarios.

\subsection{Authored physical routes}

The physics realization stages the initial six-edge tree upright at time zero. After
that boundary, it permits only joint effort and ordinary dock/undock requests. Each
action waits through the release delay, removes the old weld, follows target-relative
clearance waypoints, performs a slow measured approach, and commits after all connector
criteria pass. Routes are recorded under \pathcode{examples/scenarios}; they are not
hidden in the backend.

The routes include empirical details such as a \SI{8}{mm} lateral prebias for the final
reverse-moving three-module train and a broad first clearance tolerance. These are
transparent tuning choices. They demonstrate the framework's physical execution path,
not robustness to arbitrary initial conditions or autonomous obstacle avoidance.

\subsection{Physics ownership invariant}

The end-to-end regression replaces backend root pose, twist, and wrench methods with
functions that fail if called after scenario creation. Completing all four replacements
under that guard is stronger evidence than merely observing wheel rotation: it proves
that the scenario cannot secretly teleport its roots after staging.
